\documentclass{article}

\usepackage[final]{neurips_2024}
\makeatletter
\renewcommand{\@noticestring}{ITCS 6169/8169 Computer Vision -- Assignment 1 Report.}
\makeatother

\usepackage[utf8]{inputenc}
\usepackage[T1]{fontenc}
\usepackage{hyperref}
\usepackage{url}
\usepackage{booktabs}
\usepackage{amsfonts}
\usepackage{nicefrac}
\usepackage{microtype}
\usepackage{xcolor}
\usepackage{makecell}
\usepackage{graphicx}
\usepackage{amsmath}

% Draft marker: everything in red is a placeholder to fill in once experiments are done.
% Tighter section and float spacing to fit the two-page limit.
\usepackage{titlesec}
\titleformat{\section}{\large\bfseries}{\thesection}{1em}{}
\titlespacing*{\section}{0pt}{6pt plus 2pt minus 2pt}{3pt plus 1pt}
\setlength{\textfloatsep}{6pt plus 2pt minus 2pt}
\setlength{\intextsep}{6pt plus 2pt minus 2pt}
\setlength{\abovecaptionskip}{3pt}
\setlength{\belowcaptionskip}{2pt}

\newcommand{\todo}[1]{\textcolor{red}{[TODO: #1]}}

\title{Assignment 1: The CNN Challenge\\[2pt] \large ITCS 6169/8169 -- Computer Vision}

\author{%
  Ba-Thinh Lam \\ Department of Computer Science, University of North Carolina at Charlotte \\ \texttt{blam5@charlotte.edu}
}

\begin{document}

\maketitle
\vspace{-0.9cm}
\begin{center}
\textbf{Code \& checkpoint:} \url{https://github.com/sasukepn1999/itcs6169-assignment1} \quad
\textbf{Test accuracy:} 93.1$\pm$0.8\% \quad
\textbf{Selected-model val.\ accuracy:} 93.7$\pm$0.5\%
\end{center}

\section{Final Result}
\label{sec:result}

\textbf{Final model} (Sec.~\ref{sec:recipe}): ImageNet-pretrained EfficientNet-B0, fine-tuned with FixMatch on unlabelled images. We use three sets of images:

\textbf{Validation set} (480 labelled images held out from the provided training set, seeded split as in the starter). Every design choice is made here, with 3 seeds per configuration; we report best-epoch accuracy and the mean of the last 10 epochs, and select on the latter (the maximum of a noisy curve is optimistic). The final model's supervised part (Exp.~9: trained on the other 1{,}920 images, without FixMatch) reaches \textbf{93.7$\pm$0.5\%}.

\textbf{Test set} (the 400 labelled test images provided with the assignment). Evaluated only for final models. Our test accuracy, \textbf{93.1$\pm$0.8\%} (seeds 93.5 / 92.0 / 93.8), comes from the final recipe trained on train+val (2{,}400 images) with FixMatch on \texttt{test2}; these models never see the test images (neither labels nor pixels). For reference: best model from scratch 81.8$\pm$1.1\%, starter 42.3\%.

\textbf{\texttt{test2}} (400 unlabelled images): only used as unlabelled FixMatch data (no labels, no accuracy). \textbf{Checkpoints} (\href{https://drive.google.com/drive/folders/1q_nvnQqREGNsdmdg47nAFXf_cUVyYREe?usp=sharing}{Google Drive}): two models with the same recipe. The \emph{clean-test} models (train+val) give the test accuracy above; the \emph{submitted} model, also trained on the labelled test set, produces the \texttt{test2} predictions and has no score (no labelled data left).

\section{Data Analysis}
\label{sec:data}

\textbf{The dataset is perfectly class-balanced} (150 training and 25 test images per class), so no imbalance handling is needed. Images are far larger than $64^2$ (median $256^2$). \textbf{Image format leaks the label} (also on test): eight classes are exactly $256\times256$, Suburb $330\times220$, and \emph{Flower is the only RGB class} and the largest, so colour alone could identify it (Sec.~\ref{sec:failure}).

\section{Baseline and Experimental Journey}
\label{sec:journey}

\textbf{Starter baseline.} TNet (one $3\times3$ conv, $4\times4$ max-pool, linear classifier; 58K parameters; $64^2$ gray; 20 epochs) reaches 49.8\% validation and 42.3\% test, reproduced exactly by our code. It overfits severely (training loss 0.12; validation loss lowest at epochs 5--7) and is weakest on indoor scenes (30--35\%).

\textbf{Exp.~1: augmentation and training length.} Added one at a time, translation (reflect-pad + random crop) helps most and nearly closes the train/val gap (final loss 1.33/1.49 vs.\ 0.15/2.18); random-resized crops (RRC) and flips help less; rotation, jitter and erasing do not. Combining the three and training longer helps most, with little gain beyond 200 epochs.

\textbf{Exp.~2: learning-rate schedule.} Per-iteration decay (step, cosine, warm-up + cosine, one-cycle) beats constant lr at 200 epochs (63.4--64.1 vs.\ 62.8\%; exponential decay, too early, 60.9\%), and cosine tolerates a larger peak lr ($5\times10^{-3}$: 65.2\%, above constant lr at 600 epochs).

\textbf{Exp.~3: resolution, revisited.} With translation scaled to $s/8$, $200^2$ is best (68.3\%) and $256^2$ is worst (61.9\%), reversing a sweep made before fixing overfitting (TNet's classifier also grows with $s$).

\textbf{Exp.~4: deeper CNNs (from scratch).} TNet's training loss stays near 1.0, so it is capacity-limited. Seven torchvision CNNs from random init (better of two peak lrs kept) reach 81.6--83.5\%, except ConvNeXt-Tiny (74.5\%). \textbf{EfficientNet-B0} (4.0M params) ties ResNet-50 (23.5M) at 83.5\% with $6\times$ fewer parameters, so we keep it; it overfits again (training loss ${\approx}0.01$).

\textbf{Exp.~6--7: self-supervised auxiliary task.} An SSL loss (rotation, SimCLR, SimSiam or MAE-style masked modelling) on the training images does not reduce overfitting (rotation: within noise at $2\times$ cost; others: $-2$ to $-10$), nor on the unlabelled validation images (80.5--82.5\% test vs.\ 81.8\%).

\textbf{Exp.~8: more labels and FixMatch.} Training on all 2{,}400 train+val images raises test accuracy to 84.8$\pm$1.2\% (+3.1); FixMatch on the unlabelled test images (threshold 0.95, $\lambda_u=1$, untuned) then gives 86.9$\pm$1.4\% (+2.1, transductive, ${\approx}1.5$ std). Its pseudo-labels end up covering 98\% of the images, but their accuracy drops from ${\approx}95\%$ to 86--89\% (confirmation bias).

\textbf{Exp.~9: ImageNet initialisation.} Starting from ImageNet-1k weights (all layers fine-tuned) is the largest single gain: \textbf{93.7$\pm$0.5\%} at peak lr $10^{-4}$ vs.\ 83.5\% from scratch (+10.2); larger lrs are worse (92.5 / 90.6\%).

\begin{table}[h]
\centering
\small
\caption{Experimental journey (val.\ accuracy \%, mean$\pm$std, 3 seeds); each row builds on the one above. \emph{Best ep.}: max over epochs; \emph{Last-10}: mean of the final 10 epochs (selection metric).}
\label{tab:journey}
\begin{tabular}{lccl}
\toprule
Experiment & Best ep. & Last-10 & Observation \\
\midrule
Starter TNet ($64^2$, gray, 20 ep.) & 49.6$\pm$0.8 & 46.9$\pm$1.4 & Overfits after epoch ${\sim}7$ \\
+ translation (Exp.~1) & 54.0$\pm$1.6 & 51.0$\pm$1.2 & Rotation/jitter/erasing: no gain \\
+ flip + RRC, 100 epochs (Exp.~1b) & 63.7$\pm$1.5 & 60.5$\pm$1.1 & Translation alone: 58.4 (overfits again) \\
+ cosine lr $5\times10^{-3}$, 200 ep.\ (Exp.~2) & 66.3$\pm$1.5 & 65.2$\pm$0.9 & Constant lr, 600 ep.: 64.3 \\
+ input $200^2$ (Exp.~3) & 69.7$\pm$1.4 & 68.3$\pm$0.2 & $224^2$: 67.0; $256^2$: 61.9 \\
+ EfficientNet-B0 (Exp.~4) & 85.0$\pm$0.2 & 83.5$\pm$0.5 & ResNet-50 83.5; ConvNeXt-T 74.5 \\
\quad + rotation SSL, $\lambda=1$ (Exp.~6) & 85.8$\pm$1.0 & 84.4$\pm$0.9 & Within noise, $2\times$ cost \\
+ ImageNet init., lr $10^{-4}$ (Exp.~9) & 94.7$\pm$0.7 & \textbf{93.7$\pm$0.5} & lr $10^{-3}$: 92.5; $5\times10^{-3}$: 90.6 \\
\bottomrule
\end{tabular}
\vspace{-0.2cm}
\end{table}

\section{Our Secret Recipe}
\label{sec:recipe}

EfficientNet-B0 (4.0M parameters) initialised from \textbf{ImageNet-1k} (torchvision \texttt{IMAGENET1K\_V1}), all layers fine-tuned, new 16-way head; $200\times200$ grayscale replicated to 3 channels, ImageNet normalisation; random-resized crop (scale 0.5--1) + reflect-pad translation ($s/8$) + horizontal flip; Adam, peak lr $10^{-4}$, per-iteration cosine decay to 0; batch 64; 200 epochs; cross-entropy; \textbf{FixMatch} on the unlabelled \texttt{test2} images (pseudo-label threshold 0.95, $\lambda_u=1$, strong view: RandAugment(2,\,10) + Cutout; 64 + 64 images per step); final-epoch weights. \textbf{What mattered} (Table~\ref{tab:journey}): ImageNet initialisation (+10.2), a deeper CNN (+15.2 over TNet), augmentation with longer training (+15.9 for TNet), a cosine schedule (+2.4), and, on test (Exp.~8), more labels (+3.1) and FixMatch (+2.1). Grayscale inputs avoid a colour shortcut (Sec.~\ref{sec:failure}).

\section{Failure Analysis}
\label{sec:failure}

\textbf{Keeping colour for the colour images (Exp.~5).} \emph{Tried:} keep the originally-RGB images (only Flower) in colour. \emph{Why:} colour is genuine flower information. \emph{What happened:} validation accuracy fell (EfficientNet-B0 83.5$\to$80.9\%, ResNet-50 83.5$\to$82.8\%), and Flower accuracy did not improve (91.4\% gray vs.\ 90.3 / 84.9\%). Shown in gray at test time, Flower images are recognised only 25.8 / 18.3\% of the time (gray models: unchanged), and gray images are almost never called Flower (0.07 vs.\ 0.5\%): the models learned ``colour $\Rightarrow$ Flower'', which would pass our test set (Flower is again the only colour class) but fail on gray flowers. \emph{Learned:} a nuisance attribute perfectly correlated with one label must be removed or probed; accuracy alone can hide a shortcut.

\section{AI + Human}
\label{sec:ai}

\emph{AI help:} Claude Code help implement all the experiments and draft report. \emph{Human judgement:} I am mainly responsible for designing experiments and hypothesis based on what I observe on the results (e.g., loss curve) and my experience. Details in \texttt{AI\_USAGE.md}.

\end{document}
